% Einheit 1 von 5 – Bruch als Anteil (bank/brueche-dezimalzahlen/e1.jsonl, zusammenbau v0.1)
%% TODO Zweigzeile Teil 1: was hier gelernt wird (Satz in Schülersprache aus dem Einheitstitel) – Einheit 1
\einheitenkopf[e1]{Einheit 1 von 5 · Bruch als Anteil}
\zweigzeile{ab Kl. 5 · P10 oft}
\setcounter{aufgabe}{9}

%% TODO \verfahren{…}: Verfahrensüberschrift in Sachsprache fehlt in der Bank (2.3 b)
%% TODO Titel: Ich-kann-Satz fehlt in der Bank (Platzhalter: Kettenname) – Nr. 10
%% TODO Anweisung: ein Satz über der ersten Teilaufgabe fehlt in der Bank – Nr. 10
\begin{aufgabe}{Anteil bestimmen}
\begin{teile}
\teil Sind alle Teile gleich groß? \janein
\end{teile}
\swz{Welcher Anteil des Streifens ist gefärbt?}{\bruchrechteck{3}{5}}
\swz{Welcher Anteil des Kreises ist gefärbt?}{\bruchkreis{5}{8}}
\swz{Welcher Anteil ist gefärbt?}{\bruchrechteck{4}{9}}
\swz{Wie viel vom Kreis ist gefärbt? Gib den Bruch an.}{\bruchkreis{2}{3}}
\swz{Welcher Bruchteil des Streifens ist gefärbt?}{\bruchrechteck{5}{12}}
\swfrage{Was bleibt gleich, was ändert sich?}
\swa{Färbe $\frac{5}{6}$ des Streifens.}{\bruchrechteck{0}{6}}
\swz{Welcher Anteil vom Kreis sind die Sektoren A, B und D zusammen?}{\kreisdiagramm{A/90,B/90,C/90,D/45,E/45}}
\begin{teile}
\teil Bei welcher Figur ist $\frac{2}{5}$ gefärbt? Kreuze an.\\ \kreuz{P}\\ \kreuz{Q}\\ \kreuz{R}
\end{teile}
%% TODO Darstellung für schwach fehlt in der Bank (brueche-dezimalzahlen-e1-k1-s5-v1; Abschnitt „Für schwache Schüler“)
\swz{In einer Gruppe sind $9$ Kinder, $4$ davon tragen eine Brille. Welcher Anteil trägt eine Brille?}{}
\begin{teile}
\teil Markiert sind die Sektoren A und E. Welcher Anteil des Kreises ist markiert? \hfill (P10 2019 OS)\\ \kreuz{$\frac{2}{5}$}\\ \kreuz{$\frac{3}{9}$}\\ \kreuz{$\frac{2}{9}$}\\ \kreuz{$\frac{3}{5}$}
\end{teile}
\end{aufgabe}

%% TODO \verfahren{…}: Verfahrensüberschrift in Sachsprache fehlt in der Bank (2.3 b)
%% TODO Titel: Ich-kann-Satz fehlt in der Bank (Platzhalter: Kettenname) – Nr. 11
%% TODO Anweisung: ein Satz über der ersten Teilaufgabe fehlt in der Bank – Nr. 11
\begin{aufgabe}{Bruchteil einer Größe}
\begin{teile}
\teil Eine Tüte enthält $24$ Bonbons, $\frac{1}{3}$ davon sind rot. Was ist das Ganze?
\end{teile}
%% TODO Darstellung für schwach fehlt in der Bank (brueche-dezimalzahlen-e1-k2-s1-v1; Abschnitt „Für schwache Schüler“)
\swz{$\frac{1}{4}$ von $32$}{}
%% TODO Darstellung für schwach fehlt in der Bank (brueche-dezimalzahlen-e1-k2-s1-v2; Abschnitt „Für schwache Schüler“)
\swz{$\frac{1}{6}$ von $42$}{}
%% TODO Darstellung für schwach fehlt in der Bank (brueche-dezimalzahlen-e1-k2-s1-v3; Abschnitt „Für schwache Schüler“)
\swz{$\frac{1}{3}$ von $27$}{}
%% TODO Darstellung für schwach fehlt in der Bank (brueche-dezimalzahlen-e1-k2-s1-v4; Abschnitt „Für schwache Schüler“)
\swz{$\frac{1}{8}$ von $56$}{}
%% TODO Darstellung für schwach fehlt in der Bank (brueche-dezimalzahlen-e1-k2-s1-v5; Abschnitt „Für schwache Schüler“)
\swz{$\frac{1}{5}$ von $45$}{}
\swfrage{Was bleibt gleich, was ändert sich?}
%% TODO Darstellung für schwach fehlt in der Bank (brueche-dezimalzahlen-e1-k2-s2-v1; Abschnitt „Für schwache Schüler“)
\swz{$\frac{3}{4}$ von $36$}{}
%% TODO Darstellung für schwach fehlt in der Bank (brueche-dezimalzahlen-e1-k2-s3-v1; Abschnitt „Für schwache Schüler“)
\swz{$\frac{3}{5}$ von $45$ €}{}
%% TODO Darstellung für schwach fehlt in der Bank (brueche-dezimalzahlen-e1-k2-s4-v1; Abschnitt „Für schwache Schüler“)
\swz{Ein Kuchen ist in $12$ Stücke geteilt, $\frac{2}{3}$ davon sind gegessen. Wie viele Stücke sind noch da?}{}
%% TODO Darstellung für schwach fehlt in der Bank (brueche-dezimalzahlen-e1-k2-s5-v1; Abschnitt „Für schwache Schüler“)
\swz{$\frac{2}{5}$ von $1{,}5$ kg – wie viel Gramm?}{}
\swz[3]{Ein Aquarium fasst $120$ l und ist zu $\frac{3}{4}$ gefüllt. Wie viele Liter fehlen bis zum Rand? \hfill (P10 2015 OS)}{}
\end{aufgabe}

%% TODO Titel: Ich-kann-Satz fehlt in der Bank (Platzhalter: Kettenname) – Nr. 12
%% TODO Anweisung: ein Satz über der ersten Teilaufgabe fehlt in der Bank – Nr. 12
\begin{aufgabe}{Ganzes aus Bruchteil}
%% TODO Darstellung für schwach fehlt in der Bank (brueche-dezimalzahlen-e1-k3-s1-v1; Abschnitt „Für schwache Schüler“)
\swz{Ein Drittel der Kekse sind $8$ Stück. Wie viele Kekse sind es insgesamt?}{}
\end{aufgabe}

%% TODO Titel: Ich-kann-Satz fehlt in der Bank (Platzhalter: Kettenname) – Nr. 13
%% TODO Anweisung: ein Satz über der ersten Teilaufgabe fehlt in der Bank – Nr. 13
\begin{aufgabe}{Anteil bestimmen – Fehler finden, Begründen, Darstellungswechsel, Anwendung}
\swz[3]{Ben soll $\frac{2}{3}$ der Figur färben. Er färbt $2$ Kästchen. Finde den Fehler.}{\bruchrechteck{0}{12}}
\swfrage{Eine Pizza wird einmal unter $4$ Kinder und einmal unter $6$ Kinder gerecht verteilt. Wann ist ein Stück größer? Begründe.}
\swz{Schreibe als Bruch: fünf Neuntel.}{}
\swz[3]{Ein Fußballspiel dauert $90$ Minuten. Luca war $\frac{2}{3}$ der Spielzeit auf dem Platz. Wie viele Minuten saß er auf der Bank?}{}
\end{aufgabe}

\uebersichtskasten{Bruch: Der Nenner sagt, in wie viele gleich große Teile das Ganze geteilt ist; der Zähler sagt, wie viele Teile gemeint sind. \\ Figur aus 10 gleichen Feldern, 7 gefärbt: Anteil 7/10 \quad 4 von 20 Kindern: 4/20 = 1/5 \\ Bruchteil einer Menge: Ganzes durch den Nenner teilen, dann mal Zähler. \\ 2/5 von 35 €: 35 : 5 = 7 → 2 · 7 = 14 € \\ Rest zum Ganzen: Alle Teile zusammen sind 1. \quad 9/10 voll → 1/10 fehlt}
